% ===========================================================================
% Section 5: two pre-training predictors.
% Sources: paper_rewrite_plan.md §2 thesis item 3, §6 subplans "§4
% Trainability Is a Separate Axis" (trainability half) and "§3 The truncated
% regime" (captured-energy half), Part II §17 (key data blocks; the db3
% Hourly value is pinned there at 10.580 ± 0.102, which governs over the
% rebuttals' 10.587 and is used consistently here), Part II §20 (evidence
% inventory rows "Captured energy on Traffic", "Divergence mechanism",
% "Trend branch repairs conditioning"); rebuttal responses response_AC.md
% ("What the revised paper claims"), response_CAcc.md §2/§4,
% response_y8d1.md Q4.
%
% Statistics quoted exactly: Traffic-96 basis-level Spearman rho = -1.0,
% exact p = 0.017 (the floor at n = 5 bases); run-level rho = -0.98 over 25
% runs, p = 9.0e-18; adjacent-pair p = 0.0078 on M4-Hourly; rank sweep
% monotone p = 0.002 per rung; trend repair p = 0.00018; convergence
% tracking rho = +0.949. Tourism and Traffic applications of the resulting
% recipe are covered in Sec.~results; this section makes only the
% prediction claims.
% ===========================================================================

\section{Two Pre-training Predictors}
\label{sec:predictors}

Sections~\ref{sec:trunk} and \ref{sec:completeness} remove trunk width and
basis identity from the design space. What remains is truncation, and
truncation is decided by two quantities computable from training windows
alone, before any gradient step: captured energy
(\ref{eq:capture}) and effective rank (\ref{eq:effrank}), defined in
Section~\ref{sec:method_predictors}. They answer different questions, and
conflating them is what made the submitted advice unreliable: captured
energy predicts how accurate a truncated basis will be once trained;
effective rank predicts whether it trains at all.

\subsection{Captured energy predicts truncated-basis accuracy}
\label{sec:predictors:capture}

The cleanest demonstration fixes the rank and varies only the basis. On
M4-Hourly at $d_f = 16$ of $H = 48$, with ten seeds and byte-identical
parameter counts across arms:

\begin{center}
\small
\begin{tabular}{lcc}
\toprule
Basis & Captured energy & sMAPE \\
\midrule
Karhunen--Lo\`eve optimum & 0.999981 & $10.054 \pm 0.130$ \\
orthonormal DCT & 0.999964 & $10.287 \pm 0.097$ \\
band-ordered Daubechies-3 & 0.999885 & $10.580 \pm 0.102$ \\
SVD-ordered Daubechies-3 & 0.347765 & $73.946 \pm 4.480$ \\
\bottomrule
\end{tabular}
\end{center}

The ordering follows capture exactly, with no inversion, and every
adjacent pair differs at $p = 0.0078$ --- even where capture differs only
in the fifth decimal. Sweeping rank instead of basis gives the same law in
the other direction: as the retained fraction climbs from 0.125 to 1.000,
M4-Hourly sMAPE falls monotonically from 106.81 to 9.31, monotone at
$p = 0.002$ per rung.

The mechanism is spatial coverage, and it is visible in a single
period. On M4-Monthly ($H = 18$) the dyadic bands of the wavelet basis sit
at nine rows: truncating to six rows captures 0.658 of the target energy
and scores 73.380, while at nine rows capture reaches 0.996 and the score
13.850 --- level with a rank-matched DCT at 13.720. Three basis rows are
worth 59.5 sMAPE. The same mechanism resolves tiering, the submitted
version's offset schedule for distributing basis rows across stacks.
Rank tiering, which keeps the leading rows in every stack, lifts the mean
per-block captured energy of a DCT from 0.200 to 1.000 and beats offset
tiering by 0.138 sMAPE at $p = 0.037$; a scrambled wavelet, which has no
ordering to preserve, changes nothing ($p = 0.695$). What predicts when
the schedule matters is horizon length and band structure --- not
seasonality and not non-stationarity --- and the claim that tiering
itself helps does not survive.

Traffic-96 replicates the law out of distribution across five bases at
identical parameter counts (3{,}532{,}800): capture falls from 0.9714 to
0.3035 while sMAPE rises from 29.78 to 120.98 without a single inversion.
The basis-level correlation is Spearman $\rho = -1.0$, exact
$p = 0.017$, which is the smallest attainable value with $n = 5$ bases;
at the run level over all 25 runs, $\rho = -0.98$ with
$p = 9.0\times10^{-18}$.

\subsection{Effective rank predicts trainability}
\label{sec:predictors:effrank}

Accuracy under truncation is only half the story, because a basis can be
exactly expressive and still fail to train. At complete forecast rank ---
where Section~\ref{sec:completeness} guarantees that every basis is
equivalent --- the DCT and Karhunen--Lo\`eve arms fail on \emph{every}
seed under a plain linear synthesis head, while band-ordered and generic
arms never fail. Effective rank separates them completely:
1.05 for Karhunen--Lo\`eve and 1.54 for the DCT against 6.03
band-ordered and 17.86 and 17.12 for the two that never fail, tracking
the convergence rate at $\rho = +0.949$ --- where captured energy, flat
across all five arms to the fourth decimal, explains nothing. The
mechanism is concrete: M4-Hourly windows carry 0.98 of their energy in one
direction, so the head must learn one coefficient near 225{,}000 beside
forty-seven near 240.

Two devices repair this failure without adding expressive power. A paired
trend branch absorbs the dominant component first, taking the DCT arm
from 28.748 to 9.111 and Karhunen--Lo\`eve from 27.607 to 9.291 --- each
from zero of ten seeds converging to ten of ten at $p = 0.00018$, with a
control lacking the trend branch still at zero of ten. Rectifying the
output removes the failure by construction on non-negative targets,
cutting collapses from 107 of 23{,}786 runs to 0 of 136; the underlying
divergence mode is measured rather than hypothesized (diverged runs place
7.1--7.8\% of their points below zero, the fraction scoring the sMAPE
ceiling equals that fraction to every digit, and deleting exactly those
points recovers the converged score to within 0.11), with the full record
in Appendix~\ref{app:ablations:divergence}. Rectification destroys the
cancellation between signed corrections on sub-band designs, where it is
catastrophic (9.1 to 41.6 on the reimplementation comparison), so both
repairs are architecture-scoped rules.

The predictors also say when the trend branch is needed, which demotes it
from free insurance to a conditional rule keyed to low effective rank or
a dominant window mean; where the target's energy is not concentrated the
branch costs accuracy, and the resulting recipe and its out-of-sample
validation appear in Sections~\ref{sec:recipe} and \ref{sec:results}.
Honesty requires the counterpart: registered \emph{mechanism} predictions
around these quantities failed and are reported as failures in the
appendix register --- the rectifier diagnostic falsified two of its four
predictions, and band-ordered conditioning gave a reduced failure rate
rather than the predicted clean one --- while the quantitative predictions
from captured energy and effective rank have held. Neither quantity is
wavelet family. That is the sentence the submitted version needed and did
not have.

